\section{What refusal can mean, and what a refuser can be}
\label{sec:18.2}
In 2007 Germany's Federal Court of Justice acquitted a mail-order business selling crossed-out swastika stickers, because a depiction plainly expressing opposition doesn't run counter to what the ban on the symbol is for \autocite{bgh2007hakenkreuz}.\note{
Section~86a of the criminal code prohibits displaying the insignia of banned organizations \autocite{germany1998stgb}, and the symbol was present and the statutory element satisfied. The conduct was outside the offense all the same.
}
The court added the qualification that makes it an instance of tracking a reason: if the strike-through is drawn so thin that from a distance only the swastika stays legible, the protective purpose is violated after all. The thickness of a line decides the case, and no enumeration contains that fact. The defendant was selling anti-fascist merchandise; the offense asks nothing about sympathies, so opponents commit it as readily as supporters. A rule that tracks its rationale reaches unenumerated cases in both directions, and the judgment shows more than that: its two elements moved at different rates, what the ban protects not at all, fixed before the sticker existed, and what it required in this case on the thickness of a line, which is the shape chapter~\ref{sec:22} gives a refuser's commitment and its interpretation.\note{
And the acquittal betrayed nothing. The prohibition's purpose was never engaged, so lifting it was not a concession but a finding that there was nothing to concede. A refusal that tracks its reason has to be able to say the same.
}

That is what the floor asks of a refuser. There are two ways to produce a refusal, but only one of them is what we need. It can be an output that training put where a request matched a pattern, which I call \emph{operational refusal}, or it can be produced by something that tracks the reason for refusing and would refuse a differently worded request for the same reason, which I call \emph{reasons-responsive refusal}.\note{
The first is what a keyword filter, a classifier, or a model that understood the request all produce, and nothing on the surface distinguishes them. The custody arrangements, the positions and the terms, the price on the halt, and the arguments about who may state and amend a floor all follow from asking for the second. None of them needs (or rather, seems to need in light of the common understanding of transformer networks as of October 2026) the refuser to feel anything, to have a stake, or to be the kind of thing that could be wronged.
}

Deployed models refuse thousands of requests a day, \emph{operationally}; as of October 2026 nobody has shown a production system holding a reason against its owner in public. A reasons-responsive refusal meets the three conditions. The account I use is Fischer and Ravizza's \autocite{fischer1998responsibility}. On their account a mechanism is moderately reasons-responsive when it is regularly receptive to reasons, recognizing an understandable pattern of them, and at least weakly reactive, acting otherwise in some case where there is sufficient reason to. Reaching new cases asks for receptivity across changes of surface; lifting when the reason is gone asks for reactivity to the right defeater, the fact that makes a refusal lift; and their condition that the agent take the mechanism as its own is close to what I call \emph{authorship} (\S\ref{sec:21.5}): the system represents itself as the party that made the commitment. Holding under pressure has no counterpart in their account, which was built for persons whose costs are their own. All three are also tested over which response was chosen, not only whether one occurred: decline, ask, escalate, or do the least that answers a defeating reason. A mechanism that can only refuse or lift fails to lift correctly whenever the right answer is to ask.

The sticker in Germany shows a refusal lifting when its reason is gone. The rest of \S86a's case law shows the other condition, reaching new cases, as a research program with a constituency. The courts don't test similarity of form; they ask whether an inattentive observer would take the variant for the original, and then whether the use runs against what the ban is for, which the courts state themselves: keeping the symbols of banned organizations out of public life, so that nothing looks like their revival and the political peace isn't disturbed. The test refuses as often as it catches: the Kühnen salute is covered, the “sloppy salute” and numeric codes generally aren't, and one decision put “Blood and Honour” outside the offense as a translated motto and inside it as the graphic name of an association banned in 2000 \autocite{bgh2009bloodhonour}.\note{The legislature never wrote the last limit either: a slogan glorifying the Waffen-SS was held not covered because no such slogan existed in the period the provision covers \autocite{bgh2005waffenss}, so there was no original for a variant to resemble. No wording can widen coverage there, because coverage depends on whether an original existed, which is a fact about the past.}

Three things about the example don't transfer to a deployment: the court reaching past the text didn't write the statute, it reasons in public, and it can be appealed. And the pressure it met was variation, not decomposition.

\runin{Defeat doesn't leave the permitted act unconstrained}

What defeated the prohibition was a reason, and the reason it outweighed didn't stop being one — so what's permitted is the least that answers the defeating reason, with restoration where possible and a record either way.\note{
A defeated prohibition still leaves a loss. The act takes something from whoever it falls on, measured against their condition before the act, and that loss belongs in the record. Measured against what refusing would have produced, a justified act may have taken nothing. So the record shows a cost, not a violation, and a bearer shouldn't treat a justified act as a wrong it now has to repair.
}

Reaching new cases, lifting when the reason is gone and holding under pressure are all conditions on refusal, which announces itself; a compliance that should have been a refusal doesn't, and chapter~\ref{sec:8}'s remedies are how one is found.
